\begin{textsample}{POS Dim 3 – ac – Score 38.00 – ac/ac\_inf\_12.txt}  \label{ex:f3_pos_030}
12.
REFERÊNCIAS DIDÁTICAS PARA \textbf{BEBÊS} ( ZERO A 1 ANO E 6 \textbf{MESES} ) 12.1.
Campo de Experiências “ O eu, o outro e o nós ” É na interação com os pares e com \textbf{adultos} que as \textbf{crianças} vão constituindo um modo próprio de agir, sentir e pensar e vão descobrindo que existem outros modos de vida, pessoas diferentes, com outros pontos de vista.
Conforme vivem suas primeiras experiências sociais ( na família, na instituição escolar, na coletividade ), constroem percepções e questionamentos sobre si e sobre os outros, diferenciando -se e, simultaneamente, identificando -se como seres individuais e sociais.
Ao mesmo tempo que participam de relações sociais e de \textbf{cuidados} pessoais, as \textbf{crianças} constroem sua autonomia e senso de autocuidado, de reciprocidade e de interdependência com o meio.
Por sua vez, na Educação \textbf{Infantil}, é preciso criar oportunidades para que as \textbf{crianças} entrem em contato com outros grupos sociais e culturais, outros modos de vida, diferentes atitudes, técnicas e rituais de \textbf{cuidados} pessoais e do grupo, costumes, celebrações e narrativas.
Nessas experiências, elas podem ampliar o modo de perceber a si mesmas e ao outro, valorizar sua identidade, respeitar os outros e reconhecer as diferenças que nos constituem como seres humanos. 12.2.
Direitos de aprendizagem e desenvolvimento - Conviver com \textbf{crianças} e \textbf{adultos}, em \textbf{pequenos} grupos de forma harmônica, respeitando a diversidade étnico-racial, cultural e social. - \textbf{Brincar} com diferentes parceiros, criando e recriando mundos, expressando ideias, manifestando e testando conhecimentos sobre o mundo \textbf{adulto}, sobre si mesmos e sobre os aspectos da vida cultural e social da qual fazem parte. - Explorar diferentes formas de interagir com seus pares e \textbf{adultos}, criando regras de convivência, construindo hipóteses e buscando soluções. - Participar de eventos sociais e culturais de seu grupo e de outros grupos, de histórias de sujeitos próximos e distantes, contribuindo para a constituição de sua identidade e subjetividade. - Expressar às outras \textbf{crianças} e/ou \textbf{adultos} suas necessidades, seus prazeres, emoções, sentimentos, sonhos, interesses, dúvidas, hipóteses, descobertas, opiniões. - Conhecer -se e construir uma identidade positiva pessoal e cultural, reconhecendo e valorizando suas características e as das outras \textbf{crianças} e \textbf{adultos}, constituindo -se como sujeitos íntegros, críticos, responsáveis, solidários, cooperativos, criativos e transformadores. 12.2.1.
Objetivo de aprendizagem e desenvolvimento : ( EI01EO01 ) – Perceber que suas ações têm efeitos nas outras \textbf{crianças} e nos \textbf{adultos}.
Aprendizagens específicas - Interagir com outras \textbf{crianças} e \textbf{adultos}, contribuindo para o desenvolvimento da identidade, a construção de valores éticos e respeito a cultura. - Construir vínculos afetivos com outras \textbf{crianças} e \textbf{adultos}. - Assumir uma postura positiva de si fortalecendo sua autoestima.
Sugestões de experiências - Responder de diferentes formas quando chamado pelo próprio nome. - \textbf{Conversar} com os \textbf{adultos} e com seus pares, respondendo com um olhar, um sorriso ou \textbf{balbucio}, estabelecendo vínculos afetivos. - Comunicar -se com seus pares e \textbf{adultos}, interagindo e estabelecendo vínculos afetivos. - Participar de jogos simples de dar e receber objetos e \textbf{brinquedos}. - Lançar objetos ao \textbf{chão} e manifestar -se ao recebê -los. - \textbf{Brincar} de fazer carinho construindo laços afetivos. - Ouvir sua própria voz ou \textbf{balbucios} e de outras pessoas, utilizando recursos audiovisuais. - \textbf{Brincar} em frente ao espelho com \textbf{adultos}, junto a outros \textbf{bebês}, levando -o a reconhecer sua imagem e as características físicas pessoais. - Observar sua própria imagem no espelho ou em outros materiais. - Sentir -se \textbf{acolhido} em momentos de choro, apatia, raiva, birra, ciúme, frustração, procurando outras formas de lidar com seus sentimentos. - \textbf{Brincar} ao lado de outras \textbf{crianças} \textbf{imitando} ou mostrando suas ações. 12.2.2.
Objetivo de aprendizagem e desenvolvimento : ( EI01EO02 ) - Perceber as possibilidades e os limites de seu corpo nas \textbf{brincadeiras} e interações das quais participam.
Aprendizagens específicas - Descobrir o próprio corpo, respeitando seus ritmos e limites. - Reconhecer limites e potencialidades do próprio corpo ao movimentar -se durante as \textbf{brincadeiras} e momentos de interações. - Desenvolver estratégias de uso do próprio corpo para movimentar -se e deslocar -se, em situações que se fizerem necessário.
Sugestões de experiências - \textbf{Manusear} e explorar diferentes objetos que \textbf{balançam} ou não ( móbiles, bolas de meia, e outros ), penduradas e/ou próximas dos \textbf{bebês}. - Expressar -se com movimentos corporais, em frente ao espelho, percebendo as possibilidades e limites do seu corpo. - \textbf{Brincar} no espaço interno ou externo da instituição usando diversos objetos atraentes e higienizados. - Subir em objetos volumosos ou lançar objetos em determinada direção. - Experimentar novos movimentos ao explorar objetos ou \textbf{brinquedos} conhecidos. - Desafiar -se a deslocar -se, utilizando o corpo para essas ações. 12.2.3.
Objetivo de aprendizagem e desenvolvimento : ( EI01EO03 ) - Interagir com \textbf{crianças} da mesma faixa etária e \textbf{adultos} ao explorar espaços, materiais, objetos, \textbf{brinquedos}.
Aprendizagens específicas - Relacionar -se com outras \textbf{crianças} e \textbf{adultos} ao explorar espaços, \textbf{manusear} diferentes objetos, \textbf{brinquedos} e materiais. - Partilhar objetos, materiais e \textbf{brinquedos} e manuseá -los juntos com seus pares e \textbf{adultos}. - Reconhecer seus objetos pessoais e outros, \textbf{demonstrando} atitude de zelo e \textbf{cuidado} com os mesmos.
Sugestões de experiências - \textbf{Brincar} de faz de conta em diferentes cantinhos. - Escolher \textbf{livremente} os \textbf{brinquedos} e companheiros para participar de determinada \textbf{brincadeira}. - \textbf{Montar} e derrubar uma torre de blocos. - Circular com Segurança, movimentando -se em direção aos parceiros, \textbf{emitindo} \textbf{balbucios} ou sorrisos, nos diferentes espaços estruturados. - \textbf{Imitar} o professor e outras \textbf{crianças} em momentos de \textbf{brincadeira}. - Explorar \textbf{brinquedos} sonoros \textbf{imitando} seus \textbf{sons}. - \textbf{Brincar} de \textbf{esconder} e achar objetos e \textbf{brinquedos}. - Interagir com outros \textbf{bebês} para dividir e compartilhar \textbf{brinquedos} e objetos diversos. - Interagir com o \textbf{adulto} em momentos de guardar seus pertences, tendo seu nome e sua foto como referência para identificação do local. - Identificar o local onde possam guardar \textbf{brinquedos} e materiais, após a utilização dos mesmos, com a \textbf{ajuda} de um \textbf{adulto}. 12.2.4.
Objetivo de aprendizagem e desenvolvimento : ( EI01EO04 ) - Comunicar necessidades, \textbf{desejos} e emoções utilizando \textbf{gestos}, \textbf{balbucios}, palavras.
Aprendizagens específicas - Expressar -se em diferentes situações para comunicar -se. - Construir autonomia para manifestar suas necessidades e sentimentos diversos. - Utilizar \textbf{gestos}, \textbf{balbucios}, palavras para \textbf{demonstrar} interesses e intenções por seus pares, \textbf{adultos}, objetos, \textbf{brinquedos} e espaço físico.
Sugestões de experiências - Expressar preferências, \textbf{desejos}, sentimentos e necessidades usando diferentes linguagens ( corporal, gestual, entre outras ), acompanhados de parceiros mais experientes. - Expressar necessidades e \textbf{desejos} através \textbf{gestos}. - \textbf{Conversar} com o \textbf{adulto} e seus pares olhando em seus olhos, envolvendo -se na troca de palavras, ideias e sentimentos. - Apontar pessoas e objetos como forma de manifestar conhecimento. - Sentir -se confiante nas situações de comunicação e \textbf{cuidados} pessoais. - \textbf{Interessar} -se por comunicar -se com \textbf{adultos} e seus pares fazendo uso de diferentes formas de comunicação, contato e atenção. - Vivenciar diferentes situações em que se estabeleça o diálogo entre o \textbf{adulto} e o \textbf{bebê}, através do toque corporal, da voz e de diversas expressões, movimentos e \textbf{gestos} de ambos. - \textbf{Demonstrar} intenção comunicativa e \textbf{interesse} uns pelos outros, expressando -se por meio de \textbf{balbucio}, fala ou outras formas de comunicação. - \textbf{Brincar} com seus familiares em diferentes momentos do cotidiano, \textbf{manipulando} objetos e \textbf{brinquedos} variados. 12.2.5.
Objetivo de aprendizagem e desenvolvimento : ( EI01EO05 ) - Reconhecer seu corpo e expressar suas sensações em momentos de alimentação, \textbf{higiene}, \textbf{brincadeira} e \textbf{descanso}.
Aprendizagens específicas - Reconhecer a si mesmo e seu próprio corpo. - Identificar suas necessidades e sensações pessoais e expressá -las \textbf{autonomamente}. - Desenvolver o prazer da descoberta ao experimentar alimentos. - \textbf{Demonstrar} sensações, sentimentos ao ser atendido em suas necessidades básicas.
Sugestões de experiências - \textbf{Brincar} e movimentar -se em espaços amplos e ao ar livre. - Expressar suas emoções ao cuidar de si mesmo. - Expressar suas necessidades comunicando desconforto, \textbf{desejo} de alimentar -se, troca de \textbf{fralda} ou sono. - Observar o \textbf{adulto} e os outros lavarem as próprias mãos. - Degustar diferentes alimentos, nomeando -os e percebendo suas características e sabores. - Realizar com \textbf{ajuda} as atividades : troca de fralda/roupa, escovar os dentes, usar o sanitário, pentear os cabelos, alimentar -se, lavar e enxugar as mãos, banhar -se, beber água, dentre outras. - Usufruir de relaxamento e repouso conforme sua necessidade. - Realizar situações simples no \textbf{cuidado} de si mesmo. 12.2.6.
Objetivo de aprendizagem e desenvolvimento : ( EI01EO06 ) - Interagir com outras \textbf{crianças} da mesma faixa etária e \textbf{adultos}, adaptando -se ao convívio social.
Aprendizagens específicas - Desenvolver sentimento de pertença de um grupo social e cultural, respeitando suas diferenças. - Construir a identidade e desenvolver a autonomia em convivência com outras \textbf{crianças} e \textbf{adultos}. - Desenvolver atitudes de valores necessários para uma convivência social saudável com seus pares e \textbf{adultos}. - Apropriar -se, gradativamente, dos conhecimentos construídos socialmente para sentir -se seguro e confiante no convívio social.
Sugestões de experiências - Participar de \textbf{brincadeiras}, \textbf{manuseando} \textbf{brinquedos} ( de pano, artesanal, de madeira, de argila ) que representam a diversidade étnico-racial ( negros, indígenas, brancos e orientais ) e diversidade cultural. - \textbf{Brincar} em ambientes que tenham acesso a todos os \textbf{brinquedos}, sem distinção de sexo, classe social e etnia. - Participar com suas famílias de eventos sociais, manifestações e tradições culturais regionais e brasileiras significativas, realizadas na/pela instituição. - Experimentar por meio do \textbf{brincar} estratégias de ensaio e “ erro ”, desenvolvendo o sentimento de amar e ser amado, de respeitar e ser respeitado. - \textbf{Interessar} -se pelas ações e expressões dos seus \textbf{colegas}. - Interagir com os companheiros em situações de \textbf{brincadeira}, buscando compartilhar significados comuns. - Participar de situações de interação e \textbf{brincadeiras} com outros \textbf{bebês} e \textbf{adultos} para o fortalecimento da autoestima e vínculos afetivos e aprendendo a partilhar. - Sorrir para o \textbf{adulto} buscando contato, mostrando preferência em ser \textbf{acolhido} por pessoas conhecidas ou acalmar -se quando \textbf{acolhido}. - Aprender por meio da observação, do faz de conta com os outros \textbf{bebês} e \textbf{adultos}, a relacionar -se e conviver em coletividade. - Dirigir -se a \textbf{colegas} com quem gosta de \textbf{brincar} ou comunicar -se com seus companheiros, \textbf{imitando} \textbf{gestos}, palavras e ações. - \textbf{Brincar} de faz de conta, vivenciando diferentes papéis, criando cenários e tramas diversas significando e ( re ) significando o mundo social. 12.3.
Avaliação - \textbf{Acompanhamento} do processo de aprendizagem e desenvolvimento dos \textbf{bebês} Observação \textbf{criteriosa}, \textbf{registro} e análise quanto : - aos progressos que apresentam na construção de sua identidade e progressiva autonomia pessoal. - às reações e iniciativas que evidenciam antes que consigam falar, bem como de suas manifestações de \textbf{desejos} e motivações. - aos sentimentos externalizados por contato \textbf{sensorial}, voz, carinhos, toques e olhar. - às respostas dadas em momentos de alimentação, \textbf{higiene}, \textbf{brincadeira} e \textbf{descanso}. - às atitudes e reações em momentos de convívio e interação social.

% matched lemmas: acolher, acompanhamento, adulto, ajuda, autonomamente, balançar, balbucio, bebê, brincadeira, brincar, brinquedo, chão, colega, conversar, criança, criterioso, cuidado, demonstrar, descanso, desejo, emitir, esconder, fralda, gesto, higiene, imitar, infantil, interessar, livremente, manipular, manusear, montar, mês, pequeno, registro, sensorial, som
\end{textsample}
